\uuid{EPhZ}
\titre{Résolution d'équation différentielle du premier ordre}
\auteur{bodin ; croizat ; sacre}
\organisation{exo7}
\licence{CC BY-NC-SA 4.0}
\datecreate{2022-03-08}
\niveau{L1}
\module{Analyse}
\chapitre{Equation différentielle}
\sousChapitre{Résolution d'équation différentielle du premier ordre}
\difficulte{2}
\theme{}
\exo7qcmid{hilisit.equations-differentielles.045}

\contenu{
\begin{question}
\qclassification{225.01}
Le graphique ci-dessous représente plusieurs solutions de l'équation différentielle \(y'+2y=b\), où \(b\) est un réel. Quelle est la valeur de \(b\) ?

\qimage{images/png/EPhZ-1.png}

\begin{answers}
	\good{\(b=-2\)}
	\bad{\(b=-1\)}
	\bad{\(b=1/2\)}
	\bad{\(b=1\)}
\end{answers}
\begin{explanations}
L'équation différentielle peut s'écrire \(y'=ay+b\) avec \(a=-2\) donc ses solutions  sont les fonctions \(x\mapsto C e^{ax}-\frac ba=C e^{-2x}+\frac b2\). De plus \(\lim_{x\to +\infty} e^{-2x}=0\) donc \(b/2\) est la limite des solutions lorsque \(x\) tend vers \(+\infty\). On lit graphiquement que cette limite vaut $-1 = b/2$ donc on en déduit que $b=-2$.
\end{explanations}
\end{question}
}
